\documentclass{book}
\usepackage{fontspec}
\usepackage{lipsum}
\usepackage{polyglossia}
\setmainlanguage{french}
\usepackage{csquotes}

\begin{document}
%mettez des notes de bas de page et des notes marginales où vous voulez, observez le résultat!
% testez le changement de langue 

\chapter{Mise en forme, mise en sens}

Bien des années plus tard, face au peloton d'exécution, le colonel Aureliano Buendia devait se rappeler ce lointain après-midi au cours duquel son père l'emmena faire connaissance avec la glace. Macondo était alors un village d'une vingtaine de maisons en glaise et en roseaux, construites au bord d'une rivière dont les eaux diaphanes roulaient sur un lit de pierres polies, blanches, énormes comme des oeufs préhistoriques. Le monde était si récent que beaucoup de choses n'avaient pas encore de nom et pour les mentionner, il fallait les montrer du doigt. Tous les ans, au mois de mars, une famille de gitans déguenillés plantait sa tente près du village et, dans un grand tintamarre de fifres et de tambourins, faisait part des nouvelles inventions. Ils commencèrent par apporter l'aimant. Un gros gitan à la barbe broussailleuse et aux mains de moineau, qui répondait au nom de Melquiades, fit en public une truculente démonstration de ce que lui-même appelait la huitième merveille des savants alchimistes de Macédoine. Il passa de maison en maison, tralnant après lui deux lingots de métal, et tout le monde fut saisi de terreur à voir les chaudrons, les poêles, les tenailles et les chaufferettes tomber tout seuls de la place où ils étaient, le bois craquer à cause des clous et des vis qui essayaient désespérément de s'en arracher, et même les objets perdus depuis longtemps apparaissaient là où on les avait le plus cherchés, et se trainaient en débandade turbulente derrière les fers magiques de Melquiades. \enquote{Les choses ont une vie bien à elles, clamait le gitan avec un accent guttural; il faut réveiller leur âme, toute la question est là.}

\begin{flushright}
Gabriel Garcia Marquez, \emph{Cent ans de solitude}, Incipit
\end{flushright}

\end{document}
